% Introduction — frames around the milestone's friction-collapse finding to
% motivate Continuous GRPO. ~1 page NeurIPS-style.

\section{Introduction}
\label{sec:intro}

We study reinforcement learning for portfolio allocation across a basket of
eight broad-market ETFs (SPY, EEM, TLT, HYG, DBC, GLD, UUP, SHY) at daily
frequency. The agent picks weights on the simplex; the env rewards the
post-cost log return $r_t = \log(1 + w_t^{\top} R_{t+1} - \lambda \|w_t -
w_{t-1}\|_1)$. Three calendar splits separate training, validation, and
testing: 2008--2020 / 2021 / 2022Q1--2026Q1, the last picked because the
stock-bond correlation breaks in 2022 and standard 60/40-style baselines
that work for the prior decade lose their hedging property
(\texttt{data/splits.py:DEFAULT\_SPLIT}).

\paragraph{The friction collapse.} Our milestone established a sharp
empirical pathology in this setting: the offline RL solution that matches
classical baselines under zero friction collapses smoothly as $\lambda$
grows. Across 27 IQL runs sweeping $(\tau \in \{0.6, 0.7, 0.8\}, \beta
\in \{1, 3, 10\}, \lambda \in \{0, 0.001, 0.005\})$, the seed-pooled mean
test Sharpe was $0.94$ at $\lambda{=}0$, $0.36$ at $\lambda{=}0.001$, and
$\bm{-1.98}$ at $\lambda{=}0.005$ (47\% mean drawdown). The classical
baselines on the same test window were $0.95$ (equal weight), $0.95$
(momentum), and $1.23$ (risk parity)
(\texttt{figures/ablation\_summary.csv}; baselines from
\texttt{runs/ablation/iql\_exp\_tau0.7\_beta3.0\_tc0/.../metrics.json}).
At even modest realistic friction the offline-RL agent does not reach the
trivial 1/N baseline; at higher friction it actively destroys capital.

This is not a hyperparameter problem. The 27-run sweep shows
$\sigma_{\text{test\_sharpe}} \le 0.02$ within each friction band ---
$(\tau, \beta)$ moves the second decimal place of the answer; $\lambda$
moves the sign. The validation Sharpe consistently peaks at training
step 1{,}000 and decays through step 100{,}000, even as the offline
training loss continues to decrease. Early stopping masks the symptom; it
does not address the cause. We read this as a distribution-shift
signature: friction-induced turnover penalties are a per-step phenomenon,
but offline-RL Q-targets bootstrap from a value network that was never
trained against the friction-aware return distribution.

\paragraph{What we propose.} Two contributions, in priority order.

(1) \emph{Continuous GRPO}, an adaptation of the group-relative policy
optimization framework \citep{shao2024deepseekmath} from discrete
language-model action spaces to a continuous portfolio simplex. The key
move is environmental, not algorithmic: in our setting the next state is
provably independent of the action (Eq.~\ref{eq:exogeneity};
\texttt{tests/test\_exogeneity.py}), so the bootstrapped value baseline
in standard advantage estimation cancels under within-group centering. We
sample $G$ candidate weight vectors at each visited state, score them
through a per-state \texttt{peek\_reward\_batch} primitive added to the
env (Sec.~\ref{sec:method:grpo}), and use the within-group $z$-scored
rewards as advantages. The result is a critic-free policy gradient whose
sample efficiency is competitive with PPO and which trains directly on
the friction-aware reward signal without the offline-RL Q-target
distribution-shift pathology.

(2) \emph{Adaptive-O2O conservatism scheduling}, a per-step modulation of
the CQL penalty during offline-to-online fine-tuning that gates pessimism
on a regime-encoder KL between offline and online observation
distributions: $\beta_t = \alpha_{\mathrm{CQL}} \cdot \sigma(\eta \cdot
\mathrm{KL}(h_{\text{off}} \| h_{\text{on}}))$
(Sec.~\ref{sec:method:o2o}). Under distribution match the schedule
relaxes pessimism and lets the online updates exploit the offline
prior; under distribution shift it saturates and prevents optimistic
extrapolation onto out-of-distribution Q-estimates. We discovered during
Phase 2 prep that the existing repository's CLI plumbing of this schedule
was broken and silently bypassed the conservatism term entirely
(\texttt{writeup/draft\_implementation\_notes.md} \S 4). Every result we
report on the adaptive-vs-naive comparison post-dates that fix; the
milestone narrative around adaptive O2O could not have been produced by
the pre-fix code. We do not retroactively rewrite the milestone but flag
this in the limitations.

\paragraph{Experimental contributions.} We run all comparisons at
$\lambda \in \{0, 0.001, 0.005\}$ with three seeds. The headline
experiment (Sec.~\ref{sec:experiments:phase2c}) is a four-way comparison
at $\lambda = 0.001$: frozen offline (no fine-tune), online-only
SAC-Dirichlet (no offline pretrain), naive fine-tune, and adaptive O2O.
A GRPO group-size ablation (Sec.~\ref{sec:experiments:phase2d}) sweeps
$G \in \{4, 8, 16\}$. We pre-register the Phase 2A offline-baseline
landscape (BC, TD3+BC, CQL, AWAC, BCQ, IQL) at $\lambda = 0.001$ as the
context against which the four-way result is read.

\paragraph{What we are honest about.} The 8-ETF universe is small,
realistic-friction reward signal is weak (daily Sharpe noise on the order
of the mean), and the chronological-split protocol gives us exactly one
test window. We treat result-space variance the experiment offers ---
seed-induced and method-induced --- as the trustworthy signal, and call
out the rest as observational, in the discussion (Sec.~\ref{sec:limitations}).
